\documentclass[11pt,a4paper]{article}
\usepackage[T1]{fontenc}
\usepackage{amsmath,amssymb,amsthm,mathtools}
\usepackage{geometry}
\usepackage[hidelinks]{hyperref}
\geometry{margin=1in}

\newtheorem{theorem}{Theorem}[section]
\newtheorem{lemma}[theorem]{Lemma}
\newtheorem{proposition}[theorem]{Proposition}
\newtheorem{corollary}[theorem]{Corollary}
\theoremstyle{definition}
\newtheorem{definition}[theorem]{Definition}
\newtheorem{example}[theorem]{Example}

\author{Part II}
\date{Recorded from the lecturer in 2026-2027.}

\title{Statistical Modeling}
\begin{document}
\maketitle
\thispagestyle{empty}
\subsection*{What is a statistical model?}

\[
\begin{array}{ccc}
\text{Data modelling} & \text{vs.} & \text{Algorithmic modelling} \\
\text{Inference} & \text{vs.} & \text{Computing}
\end{array}
\]

\noindent\textbf{Def.}
A statistical model is
\[
\{P_\theta:\theta\in\Theta\}.
\]

\noindent\textbf{Examples.}
\begin{enumerate}
    \item
    \[
    \mathbb{E}\{Y\mid X=x\}=x^{\mathsf T}\beta.
    \]

    \item
    \[
    Y=X^{\mathsf T}\beta+\varepsilon,
    \qquad
    \varepsilon\mid X\sim N(0,\sigma^2).
    \]

    \item
    \[
    Y(x)=x^{\mathsf T}\beta+\varepsilon.
    \]
\end{enumerate}

\subsection*{Probability}

\[
(\Omega,\mathcal{F},P).
\]

Random variable/vector.

\begin{itemize}
    \item CDF:
    \[
    F(x)=P(X\leq x),
    \qquad x\in\mathbb{R}.
    \]

    \item PDF/PMF:
    \[
    f(x)=\frac{dF(x)}{dx}
    \qquad\text{or}\qquad
    f(x)=P(X=x).
    \]

    \item Quantile function:
    \[
    F^{-1}(u)
    =
    \inf\{x\in\mathbb{R}:F(x)\geq u\},
    \qquad u\in[0,1].
    \]

    \item MGF:
    \[
    M(t)=\mathbb{E}\{e^{tX}\},
    \qquad
    \text{for }t\text{ in a neighbourhood of }0.
    \]
\end{itemize}

\noindent
Independence. Expectation. Variance.

\medskip

\noindent
Random
\[
X=(X_1,\ldots,X_n)\in\mathbb{R}^n.
\]
Fixed
\[
A\in\mathbb{R}^{m\times n},
\qquad
b\in\mathbb{R}^m.
\]
\[
\mathbb{E}\{AX+b\}
=
A\mathbb{E}\{X\}+b.
\]

\subsection*{Normal distribution}

\[
Z\sim N_d(\mu,\Sigma),
\qquad
\mu\in\mathbb{R}^d,
\qquad
\Sigma\in\mathbb{R}^{d\times d}
\quad\text{positive definite}.
\]
\[
f(z)
=
\frac{1}{(2\pi)^{d/2}|\Sigma|^{1/2}}
e^{-\frac12(z-\mu)^{\mathsf T}\Sigma^{-1}(z-\mu)},
\qquad
z\in\mathbb{R}^d.
\]

\noindent\textbf{Properties.}
\begin{enumerate}
    \item
    \[
    AZ+b\sim N(A\mu+b,A\Sigma A^{\mathsf T}).
    \]

    \item If
    \[
    \begin{pmatrix}
        Z_1\\
        Z_2
    \end{pmatrix}
    \]
    is normal, then
    \[
    Z_1\perp\!\!\!\perp Z_2
    \iff
    \operatorname{Cov}(Z_1,Z_2)=0.
    \]
\end{enumerate}

\noindent\textbf{Informally:}
\[
\chi_d^2
=
\underbrace{
N(0,1)^2+\cdots+N(0,1)^2
}_{d\text{ times}}.
\]
\[
t_d
=
\frac{N(0,1)}{\sqrt{\chi_d^2/d}},
\qquad
F_{d_1,d_2}
=
\frac{\chi_{d_1}^2/d_1}{\chi_{d_2}^2/d_2}.
\]

\medskip

\noindent\textbf{Def.}
\begin{itemize}
    \item We say $\mathcal{P}$ is \emph{identifiable} if
    \[
    P_{\theta_1}=P_{\theta_2}
    \quad\text{iff}\quad
    \theta_1=\theta_2.
    \]

    \item We say $\mathcal{P}$ is \emph{parametric} if
    $\Theta$ is finite-dimensional
    ($\Theta\subseteq\mathbb{R}^p$).
    Otherwise non-/semi-parametric.
\end{itemize}

\medskip

\noindent\textbf{Conditional expectation}
\[
\mathbb{E}\{Y\mid X=x\}.
\]

\noindent\textbf{Law of total prob./exp.}
\[
\mathbb{E}\{Y\}
=
\mathbb{E}\bigl[\mathbb{E}\{Y\mid X\}\bigr].
\]

\subsection*{Estimation}

\noindent\textbf{Likelihood function}
\[
L(\theta)=f(\underline{X};\theta).
\]

\noindent
i.i.d.

\medskip

\noindent\textbf{Def. MLE:}
\[
\hat{\theta}
=
\operatorname*{arg\,max}_{\theta\in\Theta}L(\theta)
\qquad
\text{is a statistic}.
\]

\begin{itemize}
    \item
    \[
    \operatorname{Bias}_{\theta}(\hat{\theta})
    =
    \mathbb{E}_{\theta}(\hat{\theta})-\theta.
    \]

    \item
    \[
    \operatorname{MSE}_{\theta}(\hat{\theta})
    =
    \mathbb{E}_{\theta}\bigl\{(\hat{\theta}-\theta)^2\bigr\}
    =
    \operatorname{Bias}_{\theta}(\hat{\theta})^2
    +
    \operatorname{Var}_{\theta}(\hat{\theta}).
    \]

    \item
    $T(X)$ is \emph{sufficient} if $X\mid T$ is independent of
    $\theta$.

    \item
    \textbf{Rao--Blackwell}
    \[
    \operatorname{MSE}_{\theta}
    \left(
        \hat{\theta}=\mathbb{E}(\tilde{\theta}\mid T)
    \right)
    \leq
    \operatorname{MSE}_{\theta}(\tilde{\theta}),
    \qquad
    \text{for all }\theta.
    \]

    \item
    $S(X)\subseteq\Theta$ is a $(1-\alpha)$-confidence set if
    \[
    P_{\theta}\{\theta\in S(X)\}
    =
    1-\alpha,
    \qquad
    \forall\,\theta\in\Theta.
    \]

    Typically via pivotal quantity $R(X,\theta)$,
    distribution independent of $\theta$.
\end{itemize}

\subsection*{Hypothesis testing}

\[
H_0:\theta\in\Theta_0
\qquad\text{vs.}\qquad
H_1:\theta\in\Theta_1
\qquad
(\Theta_0,\Theta_1\text{ disjoint}).
\]

\noindent
Test function
\[
T(X)\in\{0,1\}.
\]

\subsubsection*{Neyman--Pearson theory}

\noindent\textbf{Def.}
\begin{itemize}
    \item Power function
    \[
    \beta(\theta)=P_{\theta}(T(X)=1).
    \]

    \item
    \[
    \text{Size}
    =
    \sup_{\theta\in\Theta_0}\beta(\theta).
    \]
\end{itemize}

\[
\begin{aligned}
&\text{max.}\quad \beta(\theta)\quad\text{over }\Theta_1,\\
&\text{s.t.}\quad \text{size}\leq\alpha.
\end{aligned}
\]

\subsubsection*{N--P Lemma}

If
\[
\Theta_0=\{\theta_0\},
\qquad
\Theta_1=\{\theta_1\},
\]
solution is
\[
T(X)
=
\mathbf{1}\left\{
\boxed{\frac{L(\theta_1)}{L(\theta_0)}}
\geq c(\alpha)
\right\}.
\]

\medskip

\noindent\textbf{Generalized LR:}
\[
\Lambda(X)
=
\frac{
    \displaystyle\sup_{\theta\in\Theta_0\cup\Theta_1}L(\theta)
}{
    \displaystyle\sup_{\theta\in\Theta_0}L(\theta)
}
\qquad
\text{LR statistic}.
\]

\[
\begin{array}{ll}
\text{multinomial} & /\,\mathrm{P.}\,\chi^2\\
\text{normal} & /\,\chi^2,F.
\end{array}
\]

\subsubsection*{Duality}

If
\[
T(\theta,X)
=
\mathbf{1}\{\theta\notin S(X)\},
\qquad
S(X)=\{\theta:T(\theta,X)=0\},
\]
then $T(\theta_0,X)$ is size-$\alpha$ of
$H_0:\theta=\theta_0$ for all $\theta_0$
iff $S(X)$ is a $(1-\alpha)$-confidence set.

\begin{center}
end of lecture 1
\end{center}
\end{document}
